%% kapa: exam=1 round=30 year=2019 period=2 form=A subject=accounting subject_ko=회계학 seq=32 number=72
%% source: Q-Net 감정평가사 기출문제 (cst00309) artlSeq=5205205 file=extracted/1008/5205205/2019년 제30회 감정평가사 1차시험 문제지/2019년 제30회 감정평가사 1차시험 문제지/2019년도 제30회 감정평가사 1차 2교시 A형.hwp
%% conversion: kordoc 4.24.51 -> tools/gen1.py (auto); review=proofread-sonnet; flags=none
\begin{question}{72}
(주)감평은 정상원가계산을 사용하고 있으며, 직접노무시간을 기준으로 제조간접원가를 예정배부하고 있다. (주)감평의 20x1년도 연간 제조간접원가 예산은 ￦600,000이고, 실제 발생한 제조간접원가는 ￦650,000이다. 20x1년도 연간 예정조업도는 20,000시간이고, 실제 직접노무시간은 18,000시간이다. (주)감평은 제조간접원가 배부차이를 전액 매출원가에서 조정하고 있다. 20x1년도 제조간접원가 배부차이 조정전 매출총이익이 ￦400,000이라면, 포괄손익계산서에 인식할 매출총이익은?
\begin{choices}
  \item ￦290,000
  \item ￦360,000
  \item ￦400,000
  \item ￦450,000
  \item ￦510,000
\end{choices}
\end{question}
\answer{1}
